\chapter{被测系统与 2TD-600 参考设备的使用边界}
\section{设备与现有软件构成}
\noindent\begin{tabularx}{\textwidth}{p{31mm}Y}
\toprule
对象 & 当前可从附件确认的事实\\\midrule
被测件 DUT & 单块 IMU PCB。上传工程为 Python/PySide6 上位机；现有工艺文件识别 \texttt{ASM330LHHX}，但每块实际 PCB 的芯片型号、固件、ODR 和量程仍需现场读回确认。\\
转台 & 上传维护说明书的设备型号为 \textbf{2TD-600}，U--T 结构，具备内/外轴位置、速率及摇摆控制模式。注意软件中还存在 \texttt{2FDT600WK} 温控转台分支，它是\emph{另一型号}，不能将其参数或功能无条件套用在 2TD-600 上。\\
当前工艺 & \texttt{production\_calibration.py} 与 \texttt{production\_quality\_profile.json}：ACC 12 拟合+12 验证；GYRO 静止与多速率拟合/验证；结果为试运行候选。\\
参考资料 & 《2TD-600 双轴转台使用维护说明书》，文件代号 JC10468YS（2023-03），第 3、4、8、9、10 章及附录 A。\\\bottomrule
\end{tabularx}
\section{厂家手册中应固化到实验规程的参数}
\begin{longtable}{p{51mm}p{63mm}p{34mm}}
\toprule
项目 & 手册记载 & 测试含义\\\midrule\endhead
最大负载 & $20\,\mathrm{kg}$，\textbf{包含工装与被测产品} & 装夹前称重核查\\
台面 & $\varnothing600\,\mathrm{mm}$，平面度 $0.05\,\mathrm{mm}$ & 夹具定位及误差预算\\
角位置精度 & 不大于 $\pm3''$（约 $\pm0.000833^\circ$） & 仅对应有条件的角位置指标，不等同于整个 DUT 的参考角不确定度\\
位置重复性 & 不大于 $\pm2''$ & 作为位置相关预算资料\\
两轴正交度 & 不大于 $\pm5''$ & 需考虑轴系耦合\\
内/外轴速率范围 & 两轴均约为 $\pm0.01^\circ/\mathrm{s}$ 到 $\pm60^\circ/\mathrm{s}$ & 是能力范围，\textbf{不是速率精度保证}\\
角速率精度、角加速度 & 手册均注明“无精度要求” & 不得把指令转速或反馈速率直接当高精度计量基准\\
环境 & 工作温度 $25\pm5\,{}^\circ\mathrm{C}$，RH $\le80\%$ & 本次常温实验室控制目标为 20--30\,{}$^\circ$C\\
连续工作 & 不小于 $48\,\mathrm{h}$ & 仅设备能力，非 DUT 稳定性保证\\
远程状态接口 & 手册附录 A：RS-232 至少有 10\,ms 量级状态周期，RS-422 最高 1\,ms 的发送周期配置能力 & 状态发送周期\textbf{不等于}传感器/转台之间硬件同步精度\\\bottomrule
\end{longtable}
\noindent\textit{依据：手册正文第 1--3 页（PDF 第 4--6 页）及附录 A（PDF 第 19--23 页）。该表只是摘录关键指标，具体操作以设备随机手册及现行有效校准证书为准。}
\section{机械运动约束：对测试计划的直接影响}
\begin{warningbox}
手册第 2 页对连续旋转给出\textbf{带条件的装夹方式}：台面朝上是外框俯仰零位；内框（方位）轴连续转动时外框须保持零度并静止；另一种装夹情况下外框可连续转动而内框必须保持静止。\textbf{不可把“两个轴都支持连续转动”理解为任意姿态下都能双轴无限旋转。}
\end{warningbox}
\begin{enumerate}
\item 在空载和低速人工确认外轴俯仰、内轴方位的正向、回绕逻辑、零位、锁紧装置和实际机械可达角度。
\item 按第 7、10 章安装并紧固夹具，运动前确认锁紧装置已经松开；严格遵守人员与转台\textbf{至少 3 米}的安全距离提示，检查急停、限位、线缆活动余量。
\item 使用位置反馈值作为\emph{静态位置误差}参考；复合姿态由\emph{实际反馈位置}计算，不以控制指令角度代替实测姿态。
\item 试验路径首次必须人工审查，不自动在未知机械条件下执行 $180^\circ$、$225^\circ$、$315^\circ$ 等姿态。若存在限位/线缆风险，采用\textbf{等效且已核实安全的姿态替代}，并重新检验设计矩阵条件数。
\end{enumerate}
\section{平台坐标系与内参矩阵语义}
工程现用的候选模型是外轴 X、内轴 Z：
\begin{equation}
\mathbf R_{BP}(\alpha,\beta)=\mathbf R_x(\alpha)\mathbf R_z(\beta),\qquad
\mathbf f_P=\mathbf R_{BP}^{\mathsf T}\begin{bmatrix}0&0&\sigma\end{bmatrix}^{\mathsf T}\quad[\mathrm g].
\label{eq:refacc}
\end{equation}
其中 $\alpha$ 为外轴角度，$\beta$ 为内轴角度，$\sigma=\pm1$ 是需现场确认的静止比力符号。$P$ 为夹具/平台坐标系，$B$ 为台体基座坐标系。这一模型目前依据\emph{程序约定}，还须现场对照机械正向确认。

\begin{keynote}
\textbf{不能直接把内轴 Z 当成“另一个独立倾斜轴”。} 当外轴为 $0^\circ$、内轴绕竖直方向旋转时，静止重力方向在平台坐标中不变化；因此这种运动本身\textbf{不能独立验证双轴倾角变化}。要验证横滚/俯仰的复合角度，需按实际转台运动学计算参考重力矢量，并改变外轴角与内轴角组合。
\end{keynote}
\section{参考量有效性与结果等级}
\noindent\begin{tabularx}{\textwidth}{p{38mm}p{36mm}Y}
\toprule
量 & 当前可用情况 & 允许下的结论\\\midrule
静态角位置 & 手册给出位置精度；仍需校准证书有效性和装夹外参评估 & 可以开展静态姿态对比和重复性检验\\
陀螺仪参考角速度 & 手册\textbf{没有角速率精度要求}；DLL 反馈不自动具备计量溯源 & 可进行符号、相对重复性、候选比例因子诊断；绝对比例因子精度判定需要额外速率计量保障\\
动态绝对相位/延迟 & IMU 当前有帧时间估算、转台主机轮询时间戳 & 可以诊断相对时延；缺少同步基准时，不作为硬实时延迟认证\\\bottomrule
\end{tabularx}
